% Compile with: uplatex file.tex && dvipdfmx file.dvi
%           or: latexmk -uplatex file.tex
\documentclass[dvipdfmx]{article}
\usepackage{amsmath, amssymb}
\usepackage{tikz}

\title{scenario\_bcs\_gap\_equation}
\date{\today}

\begin{document}
\maketitle

Prepare the operators and parameters for the BCS Hamiltonian. $c$ is the electron annihilation operator with momentum $k$ and spin $\sigma$ ($u$ = up, $d$ = down); $g$ is the Bogoliubov quasiparticle. $u, v$ are the transformation coefficients, $\xi$ the kinetic energy, $\Delta$ the superconducting gap.

c as the annihilation operator for electron

\[
  c
\]

g as a fermion operator

\[
  g
\]

Define u, v as scalars

\[
  u, v
\]

Define xi, Delta as scalars

\[
  \xi, \Delta
\]

Register the mean-field BCS Hamiltonian.

$H_BCS = \xi (c'(k,u) c(k,u) + c'(-k,d) c(-k,d)) + \Delta (c'(k,u) c'(-k,d) + c(-k,d) c(k,u))$

\begin{equation}
  H_BCS = \xi (c^{\dagger}(k,u) c(k,u) + c^{\dagger}(-k,d) c(-k,d)) + \Delta (c^{\dagger}(k,u) c^{\dagger}(-k,d) + c(-k,d) c(k,u))
  \label{eq:E1}
\end{equation}

The equation of motion shows how c(k,u) is coupled to H\_BCS.

Heisenberg equation of E1 for c(k,u)

\begin{equation}
  \frac{d c(k,u)}{d t} = - \Delta \mathit{i} {c}_{-k,d}^{\dagger} - \xi \mathit{i} {c}_{ku}
  \label{eq:E2}
\end{equation}

Define the Bogoliubov transformation as all four Nambu-spinor components together.

$g(k,u) = u c(k,u) - v c'(-k,d)$

\begin{equation}
  g(k,u) = u c(k,u) - v c^{\dagger}(-k,d)
  \label{eq:E3}
\end{equation}

$g(-k,d) = u c(-k,d) + v c'(k,u)$

\begin{equation}
  g(-k,d) = u c(-k,d) + v c^{\dagger}(k,u)
  \label{eq:E4}
\end{equation}

$g'(k,u) = u c'(k,u) - v c(-k,d)$

\begin{equation}
  g^{\dagger}(k,u) = u c^{\dagger}(k,u) - v c(-k,d)
  \label{eq:E5}
\end{equation}

$g'(-k,d) = u c'(-k,d) + v c(k,u)$

\begin{equation}
  g^{\dagger}(-k,d) = u c^{\dagger}(-k,d) + v c(k,u)
  \label{eq:E6}
\end{equation}

Solve the four for the electron operators (the inverse transform).

Invert E3, E4, E5, E6 for c(k,u), c(-k,d), c'(k,u), c'(-k,d)

\begin{equation}
  c_{k,u} = \frac{g^{\dagger}_{-k,d} \cdot v + g_{k,u} \cdot u}{u^{2} + v^{2}}
  \label{eq:E7}
\end{equation}

\begin{equation}
  c_{-k,d} = \frac{ - g^{\dagger}_{k,u} \cdot v + g_{-k,d} \cdot u}{u^{2} + v^{2}}
  \label{eq:E8}
\end{equation}

\begin{equation}
  c^{\dagger}_{k,u} = \frac{g^{\dagger}_{k,u} \cdot u + g_{-k,d} \cdot v}{u^{2} + v^{2}}
  \label{eq:E9}
\end{equation}

\begin{equation}
  c^{\dagger}_{-k,d} = \frac{g^{\dagger}_{-k,d} \cdot u - g_{k,u} \cdot v}{u^{2} + v^{2}}
  \label{eq:E10}
\end{equation}

Substitute the inverse transform into the Hamiltonian. Choosing u, v so the pair-creation coefficient vanishes is the Bogoliubov condition.

In E1, substitute E7 for c(k,u) and E8 for c(-k,d)

\begin{equation}
  H_BCS = \left(2 \frac{v^{2} \xi}{\left( u^{2} + v^{2} \right)^{2}} + \frac{2 \Delta u v}{\left( u^{2} + v^{2} \right)^{2}}\right)  + \left( - \frac{u v \xi}{\left( u^{2} + v^{2} \right)^{2}} + \frac{ - u v \xi}{\left( u^{2} + v^{2} \right)^{2}} - \Delta \left( \frac{ - v^{2}}{\left( u^{2} + v^{2} \right)^{2}} + \frac{u^{2}}{\left( u^{2} + v^{2} \right)^{2}} \right)\right) {g}_{-k,d}^{\dagger} {g}_{ku}^{\dagger} + \left(\frac{u^{2} \xi}{\left( u^{2} + v^{2} \right)^{2}} - \frac{v^{2} \xi}{\left( u^{2} + v^{2} \right)^{2}} - \frac{2 \Delta u v}{\left( u^{2} + v^{2} \right)^{2}}\right) {g}_{-k,d}^{\dagger} {g}_{-k,d} + \left(\frac{u^{2} \xi}{\left( u^{2} + v^{2} \right)^{2}} + \frac{ - 2 \Delta u v}{\left( u^{2} + v^{2} \right)^{2}} + \frac{ - v^{2} \xi}{\left( u^{2} + v^{2} \right)^{2}}\right) {g}_{ku}^{\dagger} {g}_{ku} + \left(2 \frac{u v \xi}{\left( u^{2} + v^{2} \right)^{2}} + \Delta \left( \frac{ - v^{2}}{\left( u^{2} + v^{2} \right)^{2}} + \frac{u^{2}}{\left( u^{2} + v^{2} \right)^{2}} \right)\right) {g}_{-k,d} {g}_{ku}
  \label{eq:E11}
\end{equation}

Now head for the gap equation. To get the Cooper-pair anomalous average, first register the pair product.

c(-k,d) c(k,u)

\begin{equation}
  c(-k,d) c(k,u)
  \label{eq:E12}
\end{equation}

Move it to the quasiparticle basis.

In E12, substitute E7 for c(k,u) and E8 for c(-k,d)

\begin{equation}
  c(-k,d) c(k,u) = \frac{u v}{\left( u^{2} + v^{2} \right)^{2}}  + \frac{v^{2}}{\left( u^{2} + v^{2} \right)^{2}} {g}_{-k,d}^{\dagger} {g}_{ku}^{\dagger} + \frac{ - u v}{\left( u^{2} + v^{2} \right)^{2}} {g}_{-k,d}^{\dagger} {g}_{-k,d} + \frac{ - u v}{\left( u^{2} + v^{2} \right)^{2}} {g}_{ku}^{\dagger} {g}_{ku} + \frac{u^{2}}{\left( u^{2} + v^{2} \right)^{2}} {g}_{-k,d} {g}_{ku}
  \label{eq:E13}
\end{equation}

Take the expectation value in the quasiparticle vacuum (the BCS ground state). Operator algebra ends here; everything after this is c-numbers.

\begin{equation}
  c(-k,d) c(k,u) = \frac{u v}{\left( u^{2} + v^{2} \right)^{2}}
  \label{eq:E14}
\end{equation}

The operators are done, so the remaining scalar algebra belongs to Maxima. Cross the session boundary — a fork, not a switch: this session stays as it is and the target becomes a new file. What travels is the value on the right plus the declarations (signs included) needed to re-establish it. The operator label on the left has no meaning in a scalar CAS, so what the value is OF is recorded in the target's note.

hand E14 off to a Maxima session

\end{document}